\subsection{A canonical congruence-cover homology detector}

\begin{statement}{RH.1}{Proved}{The level-twelve relation lattice and its first two cohomology groups}
Let $q=q_{12}:\widehat B\twoheadrightarrow Q$, where
$Q=\operatorname{SL}_2(\mathbb Z/12)/\{\pm I\}$ has order $576$,
and put $N=\ker q$ and $M=N^{\rm ab}$.  We use the genuine
normalized gentle action, with $\Phi(\tau)=\tau^{\lambda_\Phi}$.
It preserves $N$ by the integral covariance of AD.1.  Thus it
gives a continuous linear representation
\[
 h_{12}:\widehat{GT}_0\longrightarrow
          \operatorname{Aut}_{\widehat{\mathbb Z}}(M),
 \qquad M\simeq\widehat{\mathbb Z}^{97}.
\]
For the deck action of $Q$ on this lattice,
\[
 H^1(Q,M)=0,\qquad H^2(Q,M)\simeq\mathbb Z/96.
\]
Here cohomology is continuous finite-group cohomology with the
displayed profinite coefficient module. The deck representation
on $M/576M$ is faithful.
\end{statement}
\begin{proof}
Write $H_2=\langle q(s)\rangle$ and $H_3=\langle q(r)\rangle$.
The kernel $N$ meets every conjugate of either finite factor
trivially.  The ordinary Bass--Serre tree of $C_2*C_3$, divided
by the discrete congruence kernel, is therefore a finite graph
with edge set $Q$ and vertex sets $Q/H_2$ and $Q/H_3$.
Its fundamental group is free of rank
$576-288-192+1=97$.  Completion identifies its homology lattice
with $M$, and gives the exact sequence of $Q$-modules
\[
 0\longrightarrow M\longrightarrow C_1\longrightarrow C_0
       \longrightarrow\widehat{\mathbb Z}\longrightarrow0,
 \quad
 C_1=\widehat{\mathbb Z}[Q],\quad
 C_0=\widehat{\mathbb Z}[Q/H_2]\oplus
                \widehat{\mathbb Z}[Q/H_3].
 \tag{RH.1.1}
\]
All terms and maps are already defined over $\mathbb Z$; their
profinite completions are exact because the underlying graph
homology groups are free.  Equivalently one may take the inverse
limit of their exact finite-coefficient graph complexes.

Put $L=\ker(C_0\to\widehat{\mathbb Z})$.  Invariants of $C_0$
are pairs of coefficients of the two orbit sums. Their
augmentation is $(a,b)\mapsto288a+192b$. Consequently
\[
 L^Q=\widehat{\mathbb Z}(2,-3).
\]
The norm edge sum maps to exactly $(2,-3)$, up to the harmless
choice of orientation of all edges.  Shapiro's lemma gives
$H^j(Q,C_1)=0$ for $j>0$.  The long exact sequence of
$0\to M\to C_1\to L\to0$ therefore gives $H^1(Q,M)=0$.
Also $H^1(Q,C_0)=0$: by Shapiro it is the sum of the groups
$\operatorname{Hom}(C_2,\widehat{\mathbb Z})$ and
$\operatorname{Hom}(C_3,\widehat{\mathbb Z})$, both zero.
The second short exact sequence in (RH.1.1) now gives
\[
 H^2(Q,M)\simeq H^1(Q,L)
  \simeq\operatorname{coker}
   \bigl((288,192):\widehat{\mathbb Z}^{,2}
                     \to\widehat{\mathbb Z}\bigr)
  \simeq\mathbb Z/96.
\]

For $g\ne1$ in $Q$, the ordinary integral character of $M$ is
\[
 \operatorname{tr}(g\mid M)
       =1-\#(Q/H_2)^g-\#(Q/H_3)^g.
 \tag{RH.1.2}
\]
A nonidentity element cannot fix a coset of both types, because
its order would divide both two and three. Hence
$96\leq97-\operatorname{tr}(g\mid M)\leq384<576$.
If its matrix were the identity modulo $576$, this strictly
positive integer would be divisible by $576$, a contradiction.
This proves the faithful finite deck action.

The program \texttt{check\_congruence\_homology.py} constructs
the actual graph, a spanning tree, its $97$ integral cycle
generators and the deck matrices of $s,r$.  It checks their
orders, the character formula for all $576$ deck transformations,
and the primitive cusp class used in RH.3.  The cohomology
conclusions follow from the exact sequences just proved; they
are not inferred from a numerical rank computation.
\end{proof}

\begin{statement}{RH.2}{Proved}{The quadratic character is detected by homology modulo 576 at multiplier one}
On $\mathcal H$ the character $\Delta$ factors through the
finite linear action
\[
 h_{12,576}:\mathcal H\longrightarrow
        \operatorname{GL}_{97}(\mathbb Z/576).
 \tag{RH.2.1}
\]
This is an explicit canonical homology level. It is not claimed
to be the smallest detecting quotient or modulus.  In
particular, the quotient of $\widehat B$ by $N'N^{576}$ has
order $576^{98}$ and suffices for this action detector; the
larger all-homomorphisms characteristic core of FP.24 is not
required here.
\end{statement}
\begin{proof}
Suppose $\Phi\in\mathcal H$ is the identity on $M/576M$.
Equivariance of its action on $M$ with its induced action
$\Phi_Q$ on $Q$ gives, modulo $576$,
$T_{\Phi_Q(g)}=T_g$ for every $g\in Q$.
Faithfulness in RH.1 implies $\Phi_Q=1$.
For a positive integer $n$, put
\[
 \Gamma_n=\widehat B/(N'N^n),\qquad
 1\longrightarrow M/nM\longrightarrow\Gamma_n
       \longrightarrow Q\longrightarrow1.
\]
The induced automorphism of $\Gamma_{576}$ is the identity
both on its kernel and on its quotient. For a lift $b$ of
$g\in Q$, the difference $\Phi(b)b^{-1}\in M/576M$
is therefore independent of the lift and is a one-cocycle.
Its class belongs to $H^1(Q,M/576M)$.

The exact sequence $0\to M\xrightarrow{n}M\to M/nM\to0$
and RH.1 identify
\[
 H^1(Q,M/nM)\simeq H^2(Q,M)[n].
\]
Under these identifications reduction from modulus $576$ to
modulus $6$ is multiplication by $576/6=96$ on $H^2(Q,M)$.
This follows by the diagram of the two exact coefficient
sequences, whose left vertical map is multiplication by $96$
and middle vertical map is the identity.  RH.1 shows that this
reduction of the cocycle class is zero. Thus the induced
automorphism of $\Gamma_6$ is conjugation by an element of
$M/6M$.

We now use the single finite representation from FP.23, rather
than its all-homomorphisms core.  Its target is the actual image
$E_0$ of $\widehat B$ in
\[
 \bigl(\mathbb F_2[Q/H_2]\oplus
             \mathbb F_3[Q/H_3]\bigr)\rtimes Q.
\]
Its restriction to $N$ has abelian image of exponent dividing
six. Hence this representation factors through $\Gamma_6$.
Let $u$ denote the image in $E_0$ of the conjugating element;
its projection to $Q$ is one.  From
$\Phi(s)=\beta s\beta^{-1}$ and $\Phi(r)=ArA^{-1}$,
we find that $\beta^{-1}u$ centralizes the image of $s$ and
$A^{-1}u$ centralizes the image of $r$.
The two proven centralizer-projection statements of FP.23 give
$q(\beta)\in H_2$ and $q(A)\in H_3$.
Both frames are pure, and neither $H_2$ nor $H_3$ meets $q(F)$
nontrivially. Consequently
\[
 q(\beta)=q(A)=1.
\]
The pure frame abelianization, or the explicit R3 formula,
then gives $\rho\equiv0\pmod3$. The signed normalization
is the one equal to $I$ modulo four, so the R4 shear is also
$I$ modulo twelve. We obtain $T=W=I$, and therefore
$\Delta(\Phi)=0$.
This proves containment of the kernel of (RH.2.1) in
$\ker\Delta$, which is exactly the asserted factorization.
The extension $\Gamma_{576}$ has kernel
$(\mathbb Z/576)^{97}$ and quotient of order $576$, proving
the displayed order.  Its invariance under the actual gentle
action is sufficient; no claim of characteristicity under
all automorphisms of $\widehat B$ is used.
\end{proof}

\begin{statement}{RH.3}{Proved}{Full-multiplier homology factorization and its exact scope}
The full character $\Delta:\widehat{GT}_0\to\mathbb F_2$
factors continuously through the genuine representation
$h_{12}$ of RH.1, and hence through $h_{12,n}$ for some positive
integer $n$.  The value $n=576$ has been proved sufficient on
$\mathcal H$ only.  This construction provides a sign line in
a finite permutation module obtained from this homology action;
it does not assert that the sign line occurs in $M$ itself.
\end{statement}
\begin{proof}
The element $\tau^{12}$ lies in $N$. Its homology class is a
primitive nonzero vector: in the finite graph of RH.1 its
cusp cycle is an embedded cycle with coefficients $1$ and
$-1$. An edge outside any spanning tree on that cycle has
coefficient $1$ or $-1$ in its fundamental-cycle coordinates.
The program also checks this assertion in the displayed basis.
The normalized action sends this class to its multiple by
$\lambda_\Phi$. Thus $h_{12}(\Phi)=1$ forces
$\lambda_\Phi=1$; primitivity in the torsion-free profinite
lattice is essential here. RH.2 now gives $\Delta(\Phi)=0$.

The continuous map from the compact group $\widehat{GT}_0$
onto its closed image $K\subseteq\operatorname{GL}_{97}
(\widehat{\mathbb Z})$ is a quotient map. The kernel
containment just proved therefore descends $\Delta$ to a
continuous character of $K$. Its kernel is open, so it
contains $K\cap\ker(\operatorname{GL}_{97}
(\widehat{\mathbb Z})\to\operatorname{GL}_{97}
(\mathbb Z/n))$ for some $n$.
This paragraph gives a topological existence argument.
LH.1 below supplies the explicit full-multiplier bound
$n_*=2^{51942}3^{14136}$ by genuine tangential specialization
and a cyclotomic conductor estimate. Neither argument computes
the smallest modulus or a full arithmetic normalization table.

For completeness the promised module is explicit once such an
$n$ is fixed. Let $K_n$ be the actual image on $M/nM$ and let
$\mathcal O_n$ be the orbit of a chosen ordered basis of
$(\mathbb Z/n)^{97}$.  The action of $K_n$ on this orbit is
free, so the permutation module $k[\mathcal O_n]$ over any
field of characteristic different from two contains the
nonzero vector
\[
 v_\Delta=\sum_{g\in K_n}(-1)^{\Delta(g)}
                   e_{g\cdot\mathbf b}.
\]
Every group element acts on its span by $(-1)^\Delta$.
At multiplier one one may take $n=576$ throughout.
This module consists of functions on finite ordered homology
bases of the congruence cover; its construction uses the
specified tower action.  Determining the occurrence of this
character in the raw homology lattice, rather than in this
permutation functor of its finite reduction, remains a distinct
question.  LI.8 proves its absence from the canonical Alexander
module of $F$ and from the homologies of its finite abelian covers.
\end{proof}
